% =====================================================================
%  Wireless Communication — Exam Notes : shared preamble
%  Lifted verbatim from Data Mining\ExamNotes\src\preamble.tex; only the
%  tier-chip denominators and the running header differ.
%  Style: bare-bones bullets, modern sans, marks as chips,
%         figures copied from source PDFs/slides (never redrawn).
% =====================================================================
\usepackage[margin=1.45cm, top=1.35cm, bottom=1.35cm]{geometry}
\usepackage{amsmath, amssymb}
\usepackage{xcolor}
\usepackage{graphicx}
\usepackage{enumitem}
\usepackage{array}
\usepackage{tabularx}
\usepackage{booktabs}
\usepackage{colortbl}
\usepackage{titlesec}
\usepackage{fancyhdr}
\usepackage[hidelinks]{hyperref}
\usepackage{tikz}
\usetikzlibrary{calc}
\usepackage{needspace}
\usepackage[T1]{fontenc}
\usepackage[default]{lato}
% beramono silently fails to load under tectonic (no mono font ends up
% embedded and \texttt falls back to CM Roman). lmtt does load.
\renewcommand{\ttdefault}{lmtt}
\usepackage{microtype}

\parindent 0pt
\setlength{\parskip}{2pt}
\linespread{1.05}
% no orphans and no widows anywhere: a heading's year-tag line and the
% opening lines of an answer must not be cut off from each other
\clubpenalty=10000
\widowpenalty=10000
\displaywidowpenalty=10000
\brokenpenalty=10000

% ---------------------------------------------------------------- palette
\definecolor{ink}{HTML}{16202A}   % body text
\definecolor{sub}{HTML}{6B7785}   % muted meta text
\definecolor{acc}{HTML}{2563A8}   % primary accent (section, headings)
\definecolor{accL}{HTML}{D6E4F5}  % accent tint
\definecolor{mkc}{HTML}{C2410C}   % marks / warnings
\definecolor{mkL}{HTML}{FDEBD9}
\definecolor{gd}{HTML}{15803D}    % verified/added
\definecolor{gdL}{HTML}{DCFCE7}
\definecolor{pin}{HTML}{6D28D9}   % pin tier
\definecolor{ruleL}{HTML}{DEE3E8}
\definecolor{rowL}{HTML}{F4F7FA}
\color{ink}

% tight lists: this is the whole look
\setlist{topsep=1pt, itemsep=0.6pt, parsep=0pt, partopsep=0pt}
\setlist[itemize,1]{leftmargin=1.15em, label=\textcolor{acc}{\textbf{\textendash}}}
\setlist[itemize,2]{leftmargin=1.05em, label=\textcolor{sub}{\textendash}}
\setlist[itemize,3]{leftmargin=1.05em, label=\textcolor{sub}{$\circ$}}
\setlist[enumerate,1]{leftmargin=1.5em, label=\textcolor{acc}{\bfseries\arabic*.}}

% ---------------------------------------------------------------- chips
% rounded pill used for marks and tier badges
\newcommand{\pill}[3]{%
  \tikz[baseline=(p.base)]{\node[rounded corners=2.2pt, inner xsep=3.4pt,
        inner ysep=1.5pt, fill=#1, text=#2, font=\bfseries\scriptsize] (p) {#3};}}
\newcommand{\pillo}[2]{%
  \tikz[baseline=(p.base)]{\node[rounded corners=2.2pt, inner xsep=3.2pt,
        inner ysep=1.3pt, draw=#1, line width=0.4pt, text=#1,
        font=\bfseries\scriptsize] (p) {#2};}}

\newcommand{\m}[1]{\pill{mkL}{mkc}{#1}\hspace{0.5pt}}          % marks chip
\newcommand{\tS}[1]{\pill{mkc}{white}{TOP\, #1/22}}            % 6+ papers
\newcommand{\tF}[1]{\pill{acc}{white}{HOT\, #1/22}}            % 3-5 papers
\newcommand{\tP}[1]{\pill{pin}{white}{PIN\, #1/22}}            % 1-2, >=6 marks
\newcommand{\added}{\pill{gdL}{gd}{verified/added}}
\newcommand{\yr}[1]{\textcolor{sub}{\footnotesize #1}}
\newcommand{\mk}[1]{\textcolor{mkc}{\bfseries#1}}
\newcommand{\gap}{\hspace{0.55em}}
% pass/fail marks for tables and inline checks: Lato has no glyph for the
% Unicode check/cross characters (renders as a tofu box), so route through
% \checkmark (amssymb, already has a Lato-independent math glyph) and this.
% Both are control words, so TeX silently eats the space after them when
% they're followed by more prose ("\checkmark all" -> "\checkmarkall") --
% \xspace fixes that without adding a space before closing punctuation.
\usepackage{xspace}
\newcommand{\xmark}{\ensuremath{\times}\xspace}
\let\origcheckmark\checkmark
\renewcommand{\checkmark}{\origcheckmark\xspace}
% inline code / file name with a soft tint
\newcommand{\code}[1]{{\setlength{\fboxsep}{1.2pt}\colorbox{rowL}{\texttt{\small #1}}}}

% ---------------------------------------------------------------- headings
\titleformat{\section}
  {\normalfont\LARGE\bfseries\color{acc}}{\thesection.}{0.35em}{}
  [\vspace{-7pt}{\color{acc}\rule{\textwidth}{1.1pt}}]
\titlespacing*{\section}{0pt}{6pt}{7pt}

% ------------------------------------------------------- keeping a topic whole
% Lower-case \needspace reserves space by injecting stretchable glue in FRONT
% of a -100 penalty, so the reserve doubles as the width of a window in which
% TeX finds that penalty cheap: it breaks there whenever the free space is
% under roughly 4.6x the reserve, which is where the half-empty pages came
% from (at 11\baselineskip the window ran to ~640pt). Capital \Needspace, from
% the same package, instead compares \pagegoal-\pagetotal against the reserve
% and breaks only when the page really is that short. No window, so the
% reserve can be honest. Every heading below uses it.
%
% Second half of the same problem: a topic band was being stranded at the foot
% of a page with its question on the next one, and the culprit was the
% FOLLOWING heading's own reserve, which put a breakpoint between the two. \T
% raises \ifbandopen and \Q / \creamq answer it with \nobreak instead of a
% reserve of their own, so band, heading and the opening lines of the answer
% can only travel together.
% The flag has to be global (a band is often followed by a tag line wrapped in
% a group) and it has to expire on its own, or a band followed by prose rather
% than a heading would rob the NEXT question of its reserve. So the first
% paragraph started after a band clears it through \everypar and puts \everypar
% back the way it found it.
\newif\ifbandopen
\newcommand{\bandon}{\global\let\ifbandopen\iftrue
  \global\everypar{\bandoff\global\everypar{}}}
\newcommand{\bandoff}{\global\let\ifbandopen\iffalse}
\bandoff
\newcommand{\holdspace}[1]{%
  \ifbandopen \nobreak \bandoff \else \Needspace{#1\baselineskip}\fi}

% Third instalment of the same problem, 2026-08-30: the reserve is a
% guess (9 or 12 lines) and the answer under the heading is usually a \sbs
% pair three or four times that tall, so the reserve was satisfied, the
% band + heading + year-tag line were set, and then the unbreakable pair
% would not fit and went overleaf -- heading at the foot of one page, its
% whole answer on the next.
%
% Guessing a bigger reserve only trades the split for wasted paper. What
% actually fixes it is removing every legal breakpoint between the band and
% the first block of the answer, so the page builder has to fall back to the
% breakpoint BEFORE the band and move the topic whole. \nobreak (penalty
% 10000) is the only value that works: an underfull page costs 100000
% whatever penalty sits there, so 9999 would still be taken, and ties go to
% the later breakpoint.
%
% The junctions are: band|heading (already \nobreak), heading|rule,
% rule|tag, inside the tag paragraph, and tag|answer. The tag line is plain
% source text rather than a macro, so the last three are closed by \qhold
% below, which is where all the awkwardness lives.
%
% There is no escape hatch, because TeX will not give one: \unpenalty is
% illegal on the main vertical list, so the glue cannot be taken back once
% laid down, and any finite penalty loses to a later breakpoint (an underfull
% page costs 100000 whatever the penalty, and ties go to the later one). A
% pair too tall to fit under its own heading on a fresh page therefore has
% nowhere to break and would print past the bottom margin. \sbsr measures
% itself and shouts @SBS-TOO-TALL at build time instead; build.py fails on
% it. The fix for one is always the same and always in the content: split it
% into two pairs at a \lead boundary.
\makeatletter
\newcommand{\qhold}{%
  \nobreak
  % the tag line must not break in the middle either, and \clubpenalty
  % alone cannot say that for a three-line one
  \global\interlinepenalty10000
  % ... and it must not come away from the answer under it. There is no
  % LaTeX hook for "after the next paragraph": para/after would do it, but
  % switching the paragraph hooks on inserts its own node ahead of \parskip
  % and hands the page builder back the very breakpoint being closed. So
  % \par is borrowed for exactly one paragraph and put back on first use.
  \global\let\qhsavedpar\par
  \gdef\par{\qhsavedpar\global\let\par\qhsavedpar
            \global\interlinepenalty0 \nobreak}%
  % Last junction, and the one that cost the most to find: a tag line is
  % written {\color{sub}\footnotesize ...}, and \color in vertical mode puts
  % a whatsit on the list. A whatsit is not discardable, so the \parskip glue
  % that follows it becomes a legal breakpoint again and the \nobreak above
  % is wasted -- the heading would keep its rule and lose its tag. Starting
  % the paragraph here instead means \color runs in horizontal mode and
  % \parskip lands directly behind the \nobreak. Only do it when a group
  % really does come next, or a following \sbsr or list would be handed an
  % empty paragraph and an extra baseline of white.
  \qh@peek}
% The lookahead has to step over a blank line first: \T is normally followed
% by one, and that \par would otherwise both fail the group test and burn the
% one-shot hook above. Swallowing it is safe because a band always ends in
% vertical mode, so the \par would have been a no-op anyway.
\newcommand{\qh@peek}{%
  \@ifnextchar\par{\qh@eatpar}{\@ifnextchar\bgroup{\leavevmode}{}}}
\long\def\qh@eatpar#1{\qh@peek}
\makeatother

% Q = exam question heading. #1 text, #2 mark chips
\newcommand{\Q}[2]{%
  \par\vspace{7pt}\holdspace{9}%
  \begingroup\interlinepenalty10000 \noindent
  {\large\bfseries\color{ink}#1}\hspace{0.45em}#2\par\endgroup
  \nobreak\vspace{1pt}%
  \nobreak\noindent{\color{ruleL}\rule{\textwidth}{0.7pt}}\par
  \nobreak\vspace{1.5pt}\qhold}

% T = topic band (tinted full-width bar). The reserve has to cover the band,
% the question heading under it, the year-tag line and a few lines of answer.
% \qhold at the end covers the case where the band is followed by prose
% rather than by a \Q: without it the band strands at the foot of the page
% with its intro line, or with nothing at all. When a \Q does follow,
% \qhold's lookahead sees a control sequence rather than a group, so it
% leaves the paragraph unstarted and \ifbandopen survives for \Q to answer.
\newcommand{\T}[1]{%
  \par\vspace{7pt}\Needspace{12\baselineskip}\noindent
  \colorbox{accL}{\makebox[\dimexpr\textwidth-2\fboxsep][l]{%
    \textcolor{acc}{\bfseries #1}}}\par\nobreak\vspace{3pt}\bandon\qhold}

% small question sub-heading inside PYQ lists
\newcommand{\creamq}[2]{%
  \par\vspace{4pt}\holdspace{6}%
  \begingroup\interlinepenalty10000
  {\bfseries\color{acc}#1}\hspace{0.4em}#2\par\endgroup
  \nobreak\vspace{1pt}\qhold}
\let\qq\creamq

% ---------------------------------------------------------------- question box
% \begin{asked}{<year / marks tags>} ... \end{asked}
%   The exam question reproduced as the paper prints it, data table and all,
%   so a problem can be attempted from this document alone. Optional first
%   argument relabels the band for problems that are not PYQs.
%   Breakable across pages on purpose (a list, not a minipage): several
%   statements carry a full data table and will not fit in the space left.
\newenvironment{asked}[2][QUESTION AS PRINTED]%
  {\par\vspace{4pt}\Needspace{4\baselineskip}\noindent
   \colorbox{rowL}{\makebox[\dimexpr\textwidth-2\fboxsep][l]{%
     \textcolor{acc}{\bfseries\scriptsize #1}\hspace{0.7em}%
     \textcolor{sub}{\footnotesize #2}}}\par\nobreak\vspace{2pt}%
   \begin{list}{}{\leftmargin=9pt \rightmargin=0pt \topsep=0pt
                  \partopsep=0pt \parsep=1.5pt \itemsep=0pt}\item[]\small}
  {\end{list}\vspace{1pt}}

% ---------------------------------------------------------------- abbreviations
% \abbr{SHORT}{long}  — one row of the closing abbreviation chart
\newenvironment{abbrtable}
  {\par\vspace{2pt}\begin{tabularx}{\textwidth}{@{}L{1.9cm}X@{}L{1.9cm}X@{}}}
  {\end{tabularx}\par}
\newcommand{\abbr}[1]{\textbf{#1}}

% bold lead-in line
\newcommand{\lead}[1]{\vspace{2.5pt}\par\textbf{#1}\par\vspace{0.5pt}}

% ---------------------------------------------------------------- figures
\newcommand{\figC}[2]{\begin{center}\includegraphics[width=#2]{figs/#1}\end{center}}
% figure INSIDE a side-by-side column: \vspace{0pt} pins the minipage
% reference point to the TOP, else the image baseline shoves text down.
\newcommand{\figT}[1]{\vspace{0pt}\includegraphics[width=\linewidth]{figs/#1}}

% side-by-side columns
%   \sbsr{left width fraction}{left}{right}
%   \sbs{left}{right}  -> even split
% minipage on purpose: an unbreakable box, so a topic and its contents are
% never torn across a page. Costs some blank space at the foot of a page
% when a pair does not fit; that trade was made deliberately -- keeping a
% topic whole beats filling the page. Do NOT "fix" this with paracol or
% multicol: both split the pair and were reverted for exactly that reason.
% The pair is built into a box first so it can be measured. A pair taller
% than \sbsmax cannot fit under its own heading on any page, so the \nobreak
% that \qhold laid down between the tag line and this pair is taken back off
% the vertical list -- otherwise the topic would have no legal breakpoint at
% all and would print past the bottom margin. Anything over \sbsmax is a
% content bug: split it into two pairs at a \lead boundary.
\newlength{\sbsmax}\setlength{\sbsmax}{0.78\textheight}
\newsavebox{\sbsbox}
\newcommand{\sbsr}[3]{%
  \par
  \setbox\sbsbox=\hbox to\textwidth{%
  \begin{minipage}[t]{#1\textwidth}\vspace{0pt}\raggedright #2\end{minipage}%
  \hfill
  \begin{minipage}[t]{\dimexpr 0.97\textwidth - #1\textwidth\relax}%
    \vspace{0pt}\raggedright #3\end{minipage}}%
  \typeout{@SBS \the\dimexpr\ht\sbsbox+\dp\sbsbox\relax\space of \the\textheight}%
  \ifdim\dimexpr\ht\sbsbox+\dp\sbsbox\relax>\sbsmax
    \typeout{@SBS-TOO-TALL \the\dimexpr\ht\sbsbox+\dp\sbsbox\relax}%
  \fi
  \noindent\usebox{\sbsbox}\par}
\newcommand{\sbs}[2]{\sbsr{0.485}{#1}{#2}}

% ---------------------------------------------------------------- tables
\newcolumntype{L}[1]{>{\raggedright\arraybackslash}p{#1}}
\arrayrulecolor{ruleL}
\newcommand{\thead}[1]{\textcolor{acc}{\bfseries#1}}
\newcommand{\hr}{\vspace{3pt}{\color{ruleL}\rule{\textwidth}{0.7pt}}\vspace{2pt}\par}
% tinted header row for tabularx
\newcommand{\hrow}{\rowcolor{rowL}}

% ---------------------------------------------------------------- header
\pagestyle{fancy}
\fancyhf{}
\fancyhead[L]{\footnotesize\color{sub} Wireless Communication \textperiodcentered{} EX 715 / EX 751}
\fancyhead[R]{\footnotesize\color{sub}\thepage}
\renewcommand{\headrulewidth}{0pt}
\fancyheadoffset{0pt}
